\documentclass[a4paper,twoside]{article}
\usepackage{graphicx} % Required for inserting images
\usepackage[hidelinks]{hyperref} % geen kaders om linken
\usepackage{parskip}
\usepackage{bm} % For bolt symbols eg \bm{\nabla}
\usepackage{amsmath, amssymb, amsfonts} % For nicer math format
\usepackage{esint} % For bolletje in integrals
\usepackage{mathrsfs} % For curly arrr
\usepackage{subcaption}
\usepackage[english]{babel}
\usepackage{dsfont}
\usepackage[a4paper, margin=0.24cm, landscape]{geometry} 
% Verander 1.5cm naar bijvoorbeeld 1cm voor een extreem vol blad.
\usepackage[T1]{fontenc}
\usepackage[utf8]{inputenc}
\usepackage{lmodern}
\usepackage{multicol}

%Tikz packages
\usepackage{tikz}
\usetikzlibrary{angles,quotes, babel} % for pic (angle labels)
\usetikzlibrary{arrows.meta,decorations.markings,calc, decorations.pathmorphing, decorations.pathreplacing, positioning, shapes, shadings, patterns}
\usepackage{xcolor}
\usepackage{tikz-3dplot}
\usepackage[siunitx]{circuitikz}


\usepackage{etoolbox} % ifthen
\usepackage[outline]{contour} % glow around text
\usepackage{xfp} % higher precision (16 digits?)
\contourlength{1.1pt}

\setlength{\parindent}{0pt}
\setlength{\parskip}{0.08ex}
\setlength{\columnsep}{0.12cm}
%\linespread{0.8} 
%\pagestyle{empty}

\newcommand{\secblock}[1]{\vspace{0.08ex}\textbf{#1}\par}
\newcommand{\eqline}[1]{\(#1\)\par}


%\usepackage{geometry} %size regulating
%\geometry{
 %a4paper,
 %left=20mm,
 %right=20mm,
 %top=10mm,
 %bottom=20mm,
%}

% Source - https://tex.stackexchange.com/a/506180
% Posted by Henri Menke, modified by community. See post 'Timeline' for change history
% Retrieved 2026-05-07, License - CC BY-SA 4.0

%vptovf griffm.vpl
%vptovf griffb.vpl

%\DeclareFontFamily{U}{griff}{}
%\DeclareFontShape{U}{griff}{m}{n}{<-> s*[2.2] griffm}{}
%\DeclareFontShape{U}{griff}{b}{n}{<-> s*[2.2] griffb}{}
%\DeclareSymbolFont{griff}{U}{griff}{m}{n}
%\SetSymbolFont{griff}{bold}{U}{griff}{b}{n}
%\DeclareMathSymbol{\rcurs}{\mathalpha}{griff}{"72}
%\DeclareBoldMathCommand{\brcurs}{\rcurs}
%\newcommand*\hrcurs{\hat{\brcurs}}


\begin{document}


%\fontsize{4.7}{4.95}\selectfont Om alles heel kleini te maken

\begin{multicols*}{4}
\secblock{H1 Vector Analysis}
%Meeste staat op formularium
\secblock{H2 Electrostatics}

Coulomb Force: 
\(\bm F = \frac{1}{4\pi \epsilon_0}\frac{qQ}{r^2}\hat{r}\)

Electric field:
\(
\bm E (\bm r) = \frac{1}{4\pi \epsilon _0} \int \frac{\rho(\bm r ')}{r^2}\hat{\bm r} d\tau'
\)

Electric flux: 
\(
\Phi_E = \int_S \bm E \cdot d\bm a
\)

Closed path:
\(
\oint \bm E \cdot d \bm l = 0 \implies \bm\nabla
 \times \bm E = \bm 0
\)

Electric potential:
\(V(\bm r) =-\int_0^r \bm E \cdot d \bm l \implies \bm E = - \nabla V \)

% Poison/Laplace equation???

Work: 
\(W = \int _a^b \bm F \cdot d \bm l = -Q \int _a^b \bm E \cdot d \bm l  \implies W = QV(\bm r)\)

or:
\[
W = \frac{\epsilon_0}{2}\int E^2d\tau
\]

%Electrostatic pressure ? (2.51)

Capacitance:
\[
C = \frac{Q}{V} \implies W = \frac{1}{2}CV^2
\]



\secblock{H3 Potentials}

%Laplace again (maar 3.1 niet te kennen?) (3.1)

% Method of imagees (3.2) 

% Separation oof variables (3.3)

% Afleiding Legendre polynomials?

%Rodrigues formula:
%\[
%\bm P_l(x) = \frac{1}{2^l l! }(\frac{d}{dx})^l (x^2-1)^l
%\]
Curly r:
\[
\frac{1}{r_{curly}} = \frac{1}{r}\sum^\infty_{n=0} (\frac{r'}{r})^n P_n (\cos \alpha)
\]
Dipole contribution to potential:
\[
V_{\text{dip}} (\bm r) = \frac{1}{4 \pi \epsilon_0} \frac{\bm p \cdot \hat{\bm r}}{r^2}
\]
\[
\bm E_{dip}(r,\theta)=\frac{p}{4\pi \epsilon_0 r^3}(2\cos \theta \hat{\bm r} + \sin \theta \hat{\bm \theta})
\]



\secblock{H4 Electric Fields in Matter}

Atomic polarizability:
\[
\bm p = \alpha\bm E
\]

Torque:
\[
\bm N = q \bm d \times \bm E = \bm p \times \bm E
\]
\[
\bm F = (\bm p \cdot \bm \nabla)\bm E
\]
Bounded surface charge density:
\[ \bm P \cdot \hat{\bm n}\]
Bounded volume charge density:
\[
\rho_b = - \bm \nabla \cdot \bm P
\]
Potential Macroscopic field:
\[
V (\bm r) = \frac{1}{4 \pi \epsilon_0} \int \frac{\bm P(\bm r') \cdot \hat{\bm r_{c}}}{r^2_c} d\tau'
\]

Electric displacement integral form:
\[
\oint \bm D \cdot d\bm a = Q_{f_{enc}}
\]

Electric field in a homogeneous linear dielectric media:
\[
\bm E = \frac{1}{4\pi \epsilon _0} \frac{q}{r^2}\hat{\bm r}
\]
Bound charge in homogeneous isotropic linear dielectric:
\[
\rho _b = - \frac{\chi_e}{1+\chi_e} \rho_f
\]

Work:
\[
\frac{1}{2} \int \bm D \cdot \bm E d\tau
\]

% Capacitance????






\secblock{H5 Magnetostatics}
Magnetic forces do no work!\\
\(
\bm F_{mag} = I \int d \bm l \times \bm B
\)

Surface current:
\(
\bm K = \frac{d \bm I}{dl_\perp} = \sigma \bm v
\)

Magnetic force on the surface current
\(
\bm F_{mag} =  \int ( \bm v \times \bm B)\sigma da =  \int ( \bm K \times \bm B)da
\)

Volume current density:
\(
\bm J = \frac{d \bm I}{dl_\perp} = \rho \bm v
\)

Magnetic force on the volume current
\(
\bm F_{mag} =  \int ( \bm v \times \bm B)\rho d\tau =  \int ( \bm J \times \bm B)d\tau
\)

Continuity equation:
\(
\bm \nabla \cdot \bm J = \frac{\partial \rho}{\partial t}
\)

Biot-Savart law:
\(
\bm{B}(\bm{r})=\frac{\mu_0}{4\pi}\int\frac{I\times\hat{\bm{r}}_c}{r_c^2}d\bm l'
=\frac{\mu_0 I}{4\pi}\int\frac{dl'\times\hat{\bm{r}}_c}{r^2_c}
\)

For surface currents:
\(\bm{B}(\bm{r})=\frac{\mu_0}{4\pi}\int\frac{\bm K(\bm r')\times\hat{\bm{r}}_c}{r_c^2}d a'\)

For volume currents:
\(\bm{B}(\bm{r})=\frac{\mu_0}{4\pi}\int\frac{\bm J(\bm r')\times\hat{\bm{r}}_c}{r_c^2}d \tau '\)

Vector potential in magnetostatics:
\(\bm A = \frac{\mu_0}{4\pi}\int\frac{\bm J(\bm r')}{r_c}d \tau '\).

\(A_{dip}(\bm r) = \frac{\mu_0}{4\pi}\frac{\bm m \times \hat{\bm r}}{r^2}\)

magnetic dipole moment:
\(\bm m = I \bm a\)

\(\bm A_{dip} (\bm r) = \frac{\mu_0}{4\pi} \frac{m \sin \theta}{r^2}\hat{\bm \phi}\)

\(\bm B_{dip} (\bm r) = \bm \nabla \times \bm A_{dip} (\bm r) = \frac{\mu_0 m}{4\pi r^3}(2\cos \theta \hat{\bm r} + \sin \theta \hat{\bm \theta})\)
\secblock{H6 Magnetic Fields in Matter}
Torque:
\(\bm N = \bm m \times \bm B\)
Force nonuniform field "fringing":
\(\bm F = 2\pi \bm I \bm R \bm B \cos\theta\)
Force nonuniform field: 
\(\bm F = \bm \nabla (\bm m \cdot \bm B)\)

% Effect of a Magnetic field on atomic orbits?

Vector potential of a single magnetic dipole:
\(\bm A (\bm r) = \frac{\mu_0}{4\pi} \frac{\bm m \times \hat{\bm r}_c}{r^2_c} = \frac{\mu_0}{4\pi} \int \frac{\bm M (\bm r) \times \hat{\bm r}_c}{r^2_c} d\tau'\)

The potential of a volume current:
\(J_b = \bm \nabla \times \bm M\)

The potential of a surface current:
\(\bm K_b = \bm M \times \hat{\bm n}\)  





\secblock{H7 Electrodynamics}


\secblock{H8 Conservation Laws}

\secblock{H9 Electromagnetic Waves}

\secblock{H10 Potentials and Fields}

\secblock{H11 Radiation}

\secblock{H12 Electrodynamics and Relativity}

\secblock{Handige rekenregels}

\end{multicols*}
\end{document}
